\documentclass[11pt]{article}
\usepackage[margin=1in]{geometry}
\usepackage{amsmath,amssymb,amsthm}
\usepackage{hyperref}
\usepackage[numbers,sort&compress]{natbib}

\title{Representation and Tractability in Set Theoretic Estimation:\\
Sheaves, Rectangles, and the Coupling Obstruction}
\author{Project \texttt{thejeb}}
\date{\today}

\newtheorem{theorem}{Theorem}
\newtheorem{definition}{Definition}
\newtheorem{claim}{Claim}
\newtheorem*{conjecture}{Conjecture (outlook)}

\newcommand{\lean}[1]{\texttt{#1}}
\newcommand{\feas}{\mathrm{feas}}
\newcommand{\Xii}{\Xi}

\begin{document}
\maketitle

\begin{abstract}
The feasible set of a set theoretic estimation (STE) problem
\citep{combettes1993} lives in $\mathcal{P}(\Xii)$, and when the
hypothesis space is a product over $N$ variables, $|\Xii|$ is
exponential in $N$: representing the feasible set extensionally hits a
$2^N$ wall that is machine-checked in this project
(\lean{Ste.DynamicFrame.Model.card\_allExactStates}).  This note asks
when the feasible set can be represented \emph{locally} and glued, and
answers three pieces of the question with Lean-verified results in the
module \lean{Ste.Sheaf}.  (i) \emph{Gluing on the constraint side}: the
assignment $J \mapsto \feas(J)$ from constraint subsets to partial
feasible sets sends unions to intersections
(\lean{partialFeasibilitySet\_iUnion}), so any cover of the constraint
index recovers the global feasible set as the limit of local ones
(\lean{feasibilitySet\_eq\_iInter\_cover}); on this side there is no
obstruction.  (ii) \emph{A tractable upper bound}: if every constraint is
variable-separable (a rectangle), the feasible set is itself a rectangle
(\lean{rectangular\_feasibilitySet}) representable by one subset per
variable---$\sum$-space, not $\prod$-space---with product cardinality
(\lean{rectangular\_encard}).  (iii) \emph{The obstruction}: a single
coupling constraint, the diagonal on two bits, is provably not a
rectangle (\lean{diagonal\_not\_rectangular}).  This is the elementary
form of the Abramsky--Brandenburger contextuality obstruction
\citep{abramsky2011sheaf}, and the same shape as this project's
non-transitivity countermodel
(\lean{Ste.DynamicFrame.Counterexample}).  A clearly-labelled outlook
sketches the unproven cohomological program: $k$-consistency as higher
gluing, $\check{H}^1$ as the measure of intractability, and a
treewidth-style tractability dichotomy
\citep{freuder1982backtrack,dechter2003constraint}.
\end{abstract}

\section{The problem: the $2^N$ wall}

Combettes' set theoretic estimation \citep{combettes1993} specifies a
solution space $\Xii$, an index set $I$ of pieces of information, and a
property set $S_i \subseteq \Xii$ for each $i \in I$; the feasible
(solution) set is $S = \bigcap_{i} S_i$.  The framework is silent about
\emph{representation}: it treats $S$ as a set, not as a data structure.

The silence becomes a cost the moment $\Xii$ is a product.  In the
estimation problems this project cares about---dynamic frames over
documents, constraint grammars, hyperframe lattices---the hypothesis
space factors as $\Xii = \prod_{v \in V} A_v$ over a family of
variables (claims, mentions, cells), and $|\Xii| = \prod_v |A_v|$ is
exponential in $|V|$.  The finite solver in
\lean{Ste.DynamicFrame.FiniteSolver} makes the wall precise and
machine-checked: over $N$ ambient worlds there are exactly $2^N$
extensional states (\lean{Model.card\_allExactStates}), and crisp STE
has full subset expressivity---a single property set can carve out
\emph{any} subset of the universe
(\lean{arbitrary\_subset\_as\_single\_property})---so without structural
assumptions on the constraints, no representation smaller than $2^N$
can be complete.

The question of this note: \emph{which structural assumptions buy which
representations?}  The natural language for the question is
sheaf-theoretic, because ``represent locally and glue'' is exactly what
a sheaf condition asserts.

\section{Gluing over the constraint poset}

Fix $S : I \to \mathcal{P}(\Xii)$ and write, for a subset
$J \subseteq I$ of the constraint index,
\[
  \feas(J) \;=\; \bigcap_{i \in J} S_i
  \qquad (\text{\lean{partialFeasibilitySet}~$S$~$J$}),
\]
the set of hypotheses consistent with the information in $J$.  The
assignment $J \mapsto \feas(J)$ is a presheaf on the poset
$\mathcal{P}(I)$ (ordered by inclusion, mapping contravariantly): if
$J \subseteq K$ then $\feas(K) \subseteq \feas(J)$, which is the
information-monotonicity lemma
\lean{partialFeasibilitySet\_antitone} already verified in
\lean{Ste.Basic}.

The new, machine-checked content of \lean{Ste.Sheaf} is that this
presheaf satisfies the gluing condition exactly, with no sheafification
needed.  The engine is the Formal Concept Analysis bridge of
\lean{Ste.HyperFrame} \citep{ganter1999fca}: partial feasibility is
definitionally the lower polar of the satisfaction context
(\lean{partialFeasibilitySet\_eq\_lowerPolar}), and polars turn unions
into intersections.

\begin{theorem}[\lean{partialFeasibilitySet\_iUnion}]
For any family $\mathcal{J} : \kappa \to \mathcal{P}(I)$ of constraint
subsets,
\[
  \feas\Bigl(\bigcup_{k} \mathcal{J}_k\Bigr)
  \;=\; \bigcap_{k} \feas(\mathcal{J}_k).
\]
\end{theorem}

\begin{theorem}[\lean{feasibilitySet\_eq\_iInter\_cover}]
If $\bigcup_k \mathcal{J}_k = I$ is any cover of the constraint index,
then the global feasible set is the intersection of the local ones:
\[
  \feas(I) \;=\; \bigcap_{k} \feas(\mathcal{J}_k).
\]
\end{theorem}

In sheaf language: the functor $\feas$ sends colimits (unions) of
constraint sets to limits (intersections) of feasible sets, so
\emph{local consistency data over any cover glues uniquely to the
global feasible set}.  On the constraint side of the Galois connection
there is no obstruction whatsoever---this is a theorem, not a
heuristic.  Whatever makes STE hard does not live here.

\section{The tractable upper bound: rectangular constraints}

Where the hardness \emph{can} live is the variable side.  Take
$\Xii = \prod_{v \in V} A_v$.  Call a constraint \emph{rectangular}
(variable-separable, a cylinder in every coordinate) if it has the form
$\prod_{v} P_v$ for some family of per-variable sets
$P_v \subseteq A_v$---in Lean, $\lean{Set.univ.pi}\;P$.  A rectangular
constraint inspects each variable independently; it couples nothing.

\begin{theorem}[\lean{rectangular\_feasibilitySet}]
If every constraint is rectangular, $S_i = \prod_v P_{i,v}$, then the
feasible set is itself the rectangle whose $v$-th side is the pointwise
intersection:
\[
  \bigcap_{i} \prod_{v} P_{i,v}
  \;=\; \prod_{v} \Bigl(\bigcap_{i} P_{i,v}\Bigr).
\]
\end{theorem}

This is the representation theorem for the easy case: a fully
variable-separable STE problem is represented \emph{exactly} by one
subset per variable.  The representation lives in
$\sum_{v} |A_v|$ space---linear in the number of variables---never in
the $\prod_v |A_v|$ product space.  The $2^N$ wall does not exist for
rectangular constraint families.

The cardinality statement is checked as well, in terms of Mathlib's
extended natural cardinality $\mathrm{encard}$ (which handles infinite
sides gracefully):

\begin{theorem}[\lean{rectangular\_encard}]
For a finite variable set $V$ (\lean{Fintype}~$V$),
\[
  \Bigl|\bigcap_i \textstyle\prod_v P_{i,v}\Bigr|
  \;=\; \prod_{v \in V} \Bigl|\bigcap_i P_{i,v}\Bigr|,
\]
as extended cardinals.
\end{theorem}

So in the separable case even the \emph{size} of the feasible set is a
product of locally-computable sizes: consistency checking, counting,
and sampling all decompose per-variable.  This is the STE incarnation
of the trivial end of CSP tractability: a constraint network with no
edges is backtrack-free \citep{freuder1982backtrack}.

\section{The obstruction: coupling is not rectangular}

The lower bound is a single, minimal counterexample.  Let the
hypothesis space be two binary variables, $\Xii = (\mathrm{Fin}\,2 \to
\mathrm{Bool})$, and consider the \emph{diagonal} (coupling,
equality) constraint
\[
  \Delta \;=\; \{\, f \mid f(0) = f(1) \,\}
  \qquad (\text{\lean{diagonal}}).
\]

\begin{theorem}[\lean{diagonal\_not\_rectangular}]
There is no family $t : \mathrm{Fin}\,2 \to
\mathcal{P}(\mathrm{Bool})$ with $\Delta = \prod_v t_v$.
\end{theorem}

The proof is the two-line contextuality argument: $(0,0) \in \Delta$
and $(1,1) \in \Delta$, so any rectangle containing $\Delta$ has
$0 \in t_0$ and $1 \in t_1$, hence contains the mixed point $(0,1)$
---which violates the coupling.  The local (per-variable) data of
$\Delta$ is \emph{full}: projected to either coordinate, everything is
possible.  The constraint is invisible locally and exists only
jointly.

This is precisely the elementary form of the sheaf-theoretic
obstruction of Abramsky and Brandenburger \citep{abramsky2011sheaf}:
a family of perfectly consistent local sections (here: both values
possible at each variable) that admits no factorization through a
product---the ``global section'' data is not determined by, and not
recoverable from, the per-variable marginals.  In cohomological
terms (made precise only in the outlook below), the diagonal carries a
nonzero \v{C}ech gluing class relative to the cover of $V$ by
singletons.

It is also, structurally, the same countermodel this project already
uses for a different purpose: \lean{Ste.DynamicFrame.Counterexample}
refutes transitivity of possible-coreference
(\lean{maySame\_not\_transitive}) with a three-world model in which
pairwise-consistent local identifications fail to assemble into a
global one.  Pairwise consistency without global consistency is the
one phenomenon, seen twice.

Consequences for representation: since rectangles are closed under
intersection (\lean{rectangular\_feasibilitySet} again), no family of
rectangular constraints can ever express $\Delta$.  Coupling
constraints are therefore \emph{exactly} the constraints that force
representations beyond $\sum$-space; the $2^N$ wall of
\lean{card\_allExactStates} is built out of them.

\section{Outlook: the cohomological program (not proven)}

\emph{Everything in this section is conjecture or programme.  None of
it is machine-checked, and none of it should be cited as a result of
this project.}

\begin{conjecture}[Feasibility as a sheaf on a site over variables]
Section~2 proved gluing over the \emph{constraint} poset.  The
substantive open direction is the \emph{variable} side: assign to each
subset $W \subseteq V$ of variables the set of partial assignments on
$W$ consistent with all constraints whose support lies in $W$.  We
conjecture this is a presheaf on a Grothendieck site over
$\mathcal{P}(V)$ whose sheaf condition, relative to a cover of $V$ by
constraint supports, holds precisely for the tractable fragments; the
diagonal example shows the presheaf is not a sheaf for the singleton
cover in general.
\end{conjecture}

\begin{conjecture}[$k$-consistency is higher gluing]
Local ($k$-)consistency in the CSP sense
\citep{freuder1982backtrack,dechter2003constraint}---every consistent
assignment on $k-1$ variables extends to any $k$-th---should be
exactly the statement that sections glue along the cover of $V$ by
$k$-element subsets.  Freuder's theorem (strong $k$-consistency plus
width $< k$ implies backtrack-free search) would then become: on
covers refined enough for the treewidth, the presheaf \emph{is} a
sheaf, and a global section is constructed greedily.
\end{conjecture}

\begin{conjecture}[$\check{H}^1$ measures intractability]
Following \citet{abramsky2011sheaf}, the obstruction to a family of
compatible local sections extending to a global one should be captured
by a \v{C}ech cohomology class in $\check{H}^1$ of the support cover
with coefficients in the (relative) presheaf of partial solutions; the
class vanishes iff gluing succeeds.  The conjectured quantitative
version: the difficulty of an STE instance---the size blow-up of its
smallest exact representation---is controlled by the size of this
obstruction, with rectangular families the $\check{H}^1 = 0$ case of
Section~3 and the diagonal of Section~4 the minimal nonvanishing
class.
\end{conjecture}

\begin{conjecture}[Tractability dichotomy]
Bounded treewidth of the constraint (co)occurrence structure---or,
conjecturally equivalent in this setting, bounded obstruction in the
sense above---yields a polynomial-size glued representation
(junction-tree style: rectangles over bags, glued along separators),
and unbounded coupling reproduces the $2^N$ wall.  This would be the
STE form of the classical CSP dichotomy between bounded-width
tractability and general intractability
\citep{dechter2003constraint}.
\end{conjecture}

What stands \emph{proven} against this programme are its two
endpoints: the fully separable case is exactly linear
(\lean{rectangular\_feasibilitySet}, \lean{rectangular\_encard}), and
one coupling suffices to leave it
(\lean{diagonal\_not\_rectangular})---with the constraint-side gluing
(\lean{feasibilitySet\_eq\_iInter\_cover}) guaranteeing that the
difficulty is genuinely about variables, not about how information is
indexed.

\bibliographystyle{plainnat}
\bibliography{../../refs}

\end{document}
